\documentclass[11pt]{article}

\usepackage[a4paper,margin=1in]{geometry}
\usepackage{amsmath,amssymb,amsthm,mathtools}
\usepackage{enumitem}
\usepackage{xcolor}
\usepackage{hyperref}

\hypersetup{
  colorlinks=true,
  linkcolor=blue,
  urlcolor=blue,
  citecolor=blue
}

\newtheorem{theorem}{Theorem}
\theoremstyle{definition}
\newtheorem{hypothesis}{Hypothesis}
\theoremstyle{remark}
\newtheorem{remark}{Remark}

\DeclareMathOperator{\cone}{cone}
\DeclareMathOperator{\relint}{relint}

\title{Face-Selection Law at a Simple Polyhedral Vertex}
\author{}
\date{}

\begin{document}

\maketitle

Theorem V.16 in \texttt{FORMAL\_NEWTON\_TROPICAL.md} selects
$q^*=\min_j q_j$ over the monomials of the perturbation on the positive
orthant.  Section~15 of \texttt{FORMAL\_VERTEX\_THRESHOLD.md} gives the
weighted degree of a single monomial.  This note transports both results to a
simple vertex of a general polyhedron, which is what domain three adjudicates.

The transport is a change of coordinates, not a new theorem.  What is new is
that the hypothesis it needs is stated explicitly and measured per corpus row
rather than assumed.

\section{Setup}

Let $P$ be full-dimensional and let $v$ be a \emph{simple} vertex with inward
edge generators $u_1,\dots,u_n$.  Define
\[
  \Phi(c)=v+\sum_{i=1}^n c_i u_i,
  \qquad c\in\mathbb{R}_{\geq 0}^n.
\]
Simplicity makes the $u_i$ linearly independent, so $\Phi$ is a linear
isomorphism from the nonnegative orthant onto the tangent cone $T_vP$.  Set
\[
  D(c)=F(v)-F\bigl(\Phi(c)\bigr),
  \qquad
  R(c)=G\bigl(\Phi(c)\bigr)-G(v),
\]
and study
\[
  M(s)=\sup_{x\in P}\bigl(F(x)+sG(x)\bigr)
  \qquad\text{as }s\to0^+.
\]

\subsection{Hypotheses}

\begin{enumerate}[label=\arabic*.]
  \item $D(c)\geq0$ near $0$, with equality only at $0$.

  \item
  \[
    D(c)=D_0(c)+o\bigl(D_0(c)\bigr),
    \qquad
    D_0(c)=\sum_{i=1}^n A_i c_i^{\beta_i},
  \]
  where $A_i>0$ and $\beta_i>1$.

  \item The perturbation is a finite polynomial,
  \[
    R(c)=\sum_j\gamma_j c^{\alpha_j}.
  \]

  \item The relevant branch has positive effective perturbation and $q^*<1$.

  \item There is no cancellation among terms of equal lowest weighted degree
  on the selected face.
\end{enumerate}

Hypothesis~2 is the condition that the transport otherwise smuggles in.
Under the anisotropic dilation
\[
  \delta_\tau(c)
  =\bigl(\tau^{1/\beta_i}c_i\bigr)_{i=1}^n,
\]
it says exactly that
\[
  D_0(\delta_\tau c)=\tau D_0(c),
\]
so the $i$th edge coordinate has weighted degree $1/\beta_i$.
\texttt{base\_homogeneity} in the corpus measures the resulting exponent;
$1.0$ means that Hypothesis~2 holds for that row.

\section{Monomials Are Faces}

Monomial $j$ has support
\[
  S_j=\{i:\alpha_{ij}>0\}.
\]
For $S\subseteq\{1,\dots,n\}$, the orthant face
\[
  C_S=\{c\geq0:c_i=0\text{ for }i\notin S\}
\]
maps under $\Phi$ to
\[
  v+\cone\{u_i:i\in S\},
\]
the corresponding face of the tangent cone.  Moreover,
$c^{\alpha_j}$ is nonzero on $\relint(C_S)$ exactly when $S_j\subseteq S$.
Thus monomial support and tangent-cone face are the same data in edge
coordinates.

\section{Weighted Degree of a Monomial}

Under $\delta_\tau$, the monomial transforms according to
\[
  (\delta_\tau c)^{\alpha_j}
  =\tau^{q_j}c^{\alpha_j},
  \qquad
  q_j=\sum_{i=1}^n\frac{\alpha_{ij}}{\beta_i}.
\]
This value belongs to one monomial, hence to one minimal supporting face.  It
is not obtained by summing independently measured ray exponents; that was the
defect this note exists to correct.

\section{Restriction to a Face}

On $C_S$, only monomials with $S_j\subseteq S$ survive.  Define
\[
  q_S=\min\{q_j:S_j\subseteq S,\ \gamma_j\neq0\}.
\]
Factoring out $\tau^{q_S}$ gives
\[
  R\big|_{C_S}(\delta_\tau c)
  =\tau^{q_S}
   \left[
     W_S(c)
     +\sum_{q_j>q_S}\tau^{q_j-q_S}(\cdots)
   \right].
\]
Consequently,
\[
  R\big|_{C_S}(\delta_\tau c)
  \sim \tau^{q_S}W_S(c)
  \qquad (\tau\to0^+),
\]
provided the leading sum does not cancel.  A higher-dimensional face therefore
inherits the minimum $q$ among the monomials it supports.

\section{One-Parameter Balance}

Set $c=\delta_\tau z$ with $z\in C_S$.  Then
\[
  J_s(\tau,z)
  =-\tau D_0(z)+s\tau^{q_S}W_S(z)+o(\cdots).
\]
Writing
\[
  A=D_0(z)>0,
  \qquad B=W_S(z)>0,
  \qquad k=q_S,
\]
the stationary point of $-\tau A+s\tau^kB$ is
\[
  \tau_*
  =\left(\frac{skB}{A}\right)^{1/(1-k)}.
\]
Substitution gives
\[
  J=A\tau_*\frac{1-k}{k},
\]
which is positive for $0<k<1$.  Since
$\tau_*\asymp s^{1/(1-k)}$, the face predicts
\[
  \gamma_S=\frac{1}{1-q_S}.
\]

\paragraph{Check against measurement.}
For the tilted-simplex face $c_2=0$, with
$D_0=z^4$, $R=z$, and $k=1/4$, this gives the coefficient
\[
  \frac34\left(\frac14\right)^{1/3}=0.472470.
\]
The adjudicator measures $0.472470\,s^{4/3}$ to six decimal places at
$s=10^{-2},\ 2.5\times10^{-3},\ 6.25\times10^{-4}$.

\section{Selection}

For every monomial $j$, $q_{S_j}\leq q_j$, since $S_j\subseteq S_j$.
Conversely, every admissible face $S$ has $q_S=q_j$ for some surviving
monomial.  Quantifying over admissible faces and the monomials whose leading
behavior is admissible on both sides yields
\[
  q^*
  =\min_{S\text{ admissible}}q_S
  =\min_{j\text{ admissible}}q_j.
\]
Hence
\[
  M(s)-F(v)-sG(v)
  =\Theta\!\left(s^{1/(1-q^*)}\right),
  \qquad
  \boxed{\gamma=\frac{1}{1-q^*}}.
\]

The admissibility restriction is not cosmetic.  If the globally minimizing
monomial lies on an inadmissible face, the left-hand side is strictly greater.
The corpus row \texttt{3d\_shear\_orders\_2\_4\_6} has an inadmissible face at
$q=1.3804$, which does not affect the result only because it is the larger
value.

\section{Admissibility}

A face is admissible when some $z\in C_S$ satisfies
\[
  D_0(z)>0
  \qquad\text{and}\qquad
  W_S(z)>0.
\]
Admissibility fails when no monomial survives on the face, when the leading
coefficient is nonpositive, when the branch leaves the feasible cone, when the
leading terms cancel, or when $q_S\geq1$.  The tilted simplex's vertical edge
is geometrically feasible but perturbatively inactive: $R$ vanishes there
identically, so it is not a competing branch.

\section{Worked Case: The Tilted Simplex}

Let
\[
  P=\{(x_0,x_1):x_0\geq0,\ x_1\geq0,\ x_0+x_1\leq1\},
  \qquad v=(0,1),
\]
with edges
\[
  u_1=(1,-1),
  \qquad u_2=(0,-1).
\]
Thus
\[
  x=(c_1,1-c_1-c_2).
\]
For
\[
  F(x)=-\bigl((x_0+x_1-1)^2+x_0^4\bigr),
\]
we have $x_0+x_1-1=-c_2$ and $x_0=c_1$, hence
\[
  D(c)=c_2^2+c_1^4,
  \qquad
  \beta=(4,2).
\]
With $G(x)=x_0$, one has $R(c)=c_1$ and support $\{1\}$, so
\[
  q=\frac14,
  \qquad
  \gamma=\frac43.
\]
The ambient axis $e_0=(1,0)$ satisfies $e_{0,0}+e_{0,1}=1>0$ and therefore
does not belong to the tangent cone.  Its quadratic decay is irrelevant to the
constrained problem.  This states the \texttt{ambient\_exponent\_law}
counterexample exactly.

\section{Scope and What Is Not Proved}

The result is conditional: for a simple vertex, if edge coordinates transport
the tangent cone to an orthant, if the base has a valid weighted principal
part, if the perturbation has noncancelling polynomial leading terms, and if
Theorem~V.16's admissibility conditions hold, then
\[
  \boxed{
    \gamma=\frac{1}{1-\min_{S\text{ admissible}}q_S}
  }.
\]

The following claims are not established here.

\begin{itemize}
  \item \textbf{Hypothesis~2 for arbitrary $F$.}
  In edge coordinates a base may contain cross terms.  A cross term of lower
  weight makes the drop $\Theta(\tau^c)$ with $c<1$, so the per-edge
  $\beta_i$ cease to describe the cone's interior.  Constructed violations at
  $c=0.75$ and $c=0.5$ still predict correctly---a positive cross term only
  steepens the interior and drives the optimum onto a face---but agreement is
  not a license.  This is why \texttt{base\_homogeneity} is recorded per row
  and is never gated on.

  \item \textbf{Nonsimplicial vertices.}
  With more than $n$ active constraints, $\Phi$ is not an isomorphism from the
  orthant, and the argument fails at the first step.  The vertex probe
  currently returns the vertex value there rather than answering.
\end{itemize}

Neither limitation is a gap in the coordinate transport; both delimit where
it applies.

\end{document}
